\documentclass[11pt,a4paper]{article}
\usepackage[utf8]{inputenc}
\usepackage[spanish,es-tabla]{babel}
\usepackage{amsmath,amssymb,amsfonts,amsthm}
\usepackage{geometry}
\geometry{top=2.5cm,bottom=2.5cm,left=2.5cm,right=2.5cm}
\usepackage{graphicx}
\usepackage{xcolor}
\usepackage{tcolorbox}
\tcbuselibrary{skins,breakable}
\usepackage{booktabs}
\usepackage{array}
\usepackage{multirow}
\usepackage{hyperref}
\usepackage{fancyhdr}
\usepackage{colortbl}

% Colores institucionales y pedagógicos
\definecolor{primary}{RGB}{30, 41, 59}     % Slate oscuro
\definecolor{secondary}{RGB}{15, 118, 110}  % Teal elegante
\definecolor{accentgreen}{RGB}{16, 185, 129}% Verde esmeralda
\definecolor{accentblue}{RGB}{2, 132, 199}  % Azul cielo
\definecolor{softpeach}{RGB}{254, 226, 226} % Rosa coral
\definecolor{softgreen}{RGB}{220, 252, 231} % Verde pastel
\definecolor{softlavender}{RGB}{237, 237, 250}% Lavanda pastel
\definecolor{softorange}{RGB}{254, 243, 199}% Melocotón pastel
\definecolor{lightgray}{RGB}{241, 245, 249}

\hypersetup{
    colorlinks=true,
    linkcolor=secondary,
    urlcolor=accentblue,
    citecolor=secondary
}

\pagestyle{fancy}
\fancyhf{}
\rhead{\textcolor{secondary}{\textbf{Análisis Estructural II} | Cátedra Ing. Ulianov Cuba Valencia}}
\lhead{\textcolor{primary}{Manual Teórico Maestro: Método de la Rigidez}}
\rfoot{\thepage}
\renewcommand{\headrulewidth}{0.8pt}
\renewcommand{\footrulewidth}{0.4pt}

\newtcolorbox{conceptbox}[1]{
    colback=lightgray,
    colframe=secondary,
    fonttitle=\bfseries\color{white},
    title=#1,
    arc=4mm,
    boxrule=1pt,
    breakable
}

\newtcolorbox{formulabox}[1]{
    colback=softgreen,
    colframe=accentgreen,
    fonttitle=\bfseries\color{primary},
    title=#1,
    arc=3mm,
    boxrule=1pt,
    breakable
}

\newtcolorbox{examtip}[1]{
    colback=softorange,
    colframe=orange!80!black,
    fonttitle=\bfseries\color{primary},
    title=CONSEJO DE EXAMEN / REGLA DE ORO: #1,
    arc=3mm,
    boxrule=1.2pt,
    breakable
}

\begin{document}

% -------------------------------------------------------------
% PORTADA
% -------------------------------------------------------------
\begin{titlepage}
    \centering
    \vspace*{1cm}
    {\scshape\LARGE Universidad Continental \par}
    \vspace{0.4cm}
    {\scshape\large Facultad de Ingeniería | Carrera Profesional de Ingeniería Civil \par}
    \vspace{1.5cm}
    {\Large\bfseries ASIGNATURA: ANÁLISIS ESTRUCTURAL II \par}
    \vspace{0.3cm}
    {\large\textit{Docente de Cátedra: Ing. Ulianov Cuba Valencia} \par}
    \vspace{2cm}
    \rule{\linewidth}{1.5pt} \\[0.4cm]
    {\huge\bfseries MANUAL TEÓRICO MAESTRO Y TUTORIAL PASO A PASO: MÉTODO DE LA RIGIDEZ DIRECTA EN ARMADURAS PLANAS \par}
    \vspace{0.2cm}
    {\Large\itshape Fundamentación Física, Deducción Matricial, Partición Cinemática (Por qué se eliminan filas y columnas) y Automatización Profesional en Excel \par}
    \rule{\linewidth}{1.5pt} \\[1.5cm]
    \vfill
    {\large\textbf{Autor / Estudiante:} Ingeniería Civil -- Universidad Continental \par}
    {\large\textbf{Alcance:} Modelos Universales de hasta 6 Nodos (12 GDL) y 8 Barras \par}
    \vspace{0.8cm}
    {\large Perú \\ Semestre 2026-II \par}
\end{titlepage}

\tableofcontents
\newpage

% -------------------------------------------------------------
\section{Introducción y Fundamento Físico de las Armaduras Planas}
% -------------------------------------------------------------

Una armadura plana (reticulado o cercha) es un ensamble de elementos rectilíneos esbeltos interconectados en sus extremos mediante nudos articulados (pasadores o rótulas libres de fricción).

\subsection{Hipótesis Clásicas del Análisis Matricial de Armaduras}
\begin{enumerate}
    \item \textbf{Barras Biarticuladas:} En los extremos no se transmiten momentos flectores ($M = 0$) ni fuerzas cortantes ($V = 0$). Cada barra trabaja \textbf{exclusivamente a esfuerzo axial} (tensión o compresión pura).
    \item \textbf{Cargas Aplicadas Únicamente en los Nudos:} Si existieran cargas distribuidas sobre las barras, la estática clásica exige calcular previamente sus reacciones equivalentes en los nudos.
    \item \textbf{Comportamiento Elástico Lineal (Ley de Hooke):} Se asume que el material cumple la ley de proporcionalidad directa entre esfuerzos y deformaciones ($\sigma = E \cdot \varepsilon$), y que las deformaciones son pequeñas (teoría de primer orden / pequeñas deformaciones).
\end{enumerate}

\subsection{Deducción de la Rigidez Axial Unidimensional}
Para una barra elástica homogénea de longitud $L$, área transversal $A$ y módulo de elasticidad $E$:
\begin{equation}
    \sigma = E \cdot \varepsilon \implies \frac{F}{A} = E \left(\frac{\Delta L}{L}\right) \implies F = \left(\frac{EA}{L}\right) \Delta L
\end{equation}
Definimos la \textbf{rigidez axial de la barra} como:
\begin{equation}
    k_{axial} = \frac{EA}{L} \quad \left[\frac{\text{Kgf}}{\text{cm}}\right]
\end{equation}
La rigidez $k_{axial}$ representa la fuerza necesaria para producir un alargamiento o acortamiento unitario ($\Delta L = 1\text{ cm}$) en la dirección del eje local de la barra.

% -------------------------------------------------------------
\section{De Coordenadas Locales a Globales: Matriz de Rigidez $4 \times 4$}
% -------------------------------------------------------------

En el plano global $(X, Y)$, cada barra $m$ conecta un \textbf{Nudo Inicial ($i$)} con un \textbf{Nudo Final ($j$)}. Cada nudo posee dos grados de libertad traslacionales:
\begin{equation}
    \{D_m\}_{global} = \begin{bmatrix} u_{iX} \\ u_{iY} \\ u_{jX} \\ u_{jY} \end{bmatrix}
\end{equation}

\subsection{Cosenos Directores y Geometría}
Dados los puntos $i(X_i, Y_i)$ y $j(X_j, Y_j)$:
\begin{align}
    \Delta X &= X_j - X_i, \quad \Delta Y = Y_j - Y_i \\
    L &= \sqrt{(\Delta X)^2 + (\Delta Y)^2} \\
    C_x &= \cos\theta = \frac{\Delta X}{L}, \quad C_y = \sin\theta = \frac{\Delta Y}{L}
\end{align}
Cumpliéndose la identidad trigonométrica fundamental: $C_x^2 + C_y^2 = 1$.

\subsection{Matriz de Transformación y Rotación de Coordenadas}
La proyección del desplazamiento global sobre el eje longitudinal de la barra viene dada por la matriz de transformación $[T]$:
\begin{equation}
    [T] = \begin{bmatrix} C_x & C_y & 0 & 0 \\ 0 & 0 & C_x & C_y \end{bmatrix}
\end{equation}

Por el principio del trabajo virtual o energía de deformación, la matriz de rigidez en coordenadas globales $[k_m]$ es:
\begin{equation}
    [k_m] = [T]^T [k'_{local}] [T]
\end{equation}

Desarrollando el producto matricial obtenemos la célebre \textbf{Matriz de Rigidez Local $4 \times 4$}:
\begin{formulabox}{Matriz de Rigidez Local de Barra en Coordenadas Globales}
\begin{equation}
    [k_m] = \frac{EA}{L}
    \begin{bmatrix}
        C_x^2     & C_x C_y   & -C_x^2    & -C_x C_y  \\
        C_x C_y   & C_y^2     & -C_x C_y  & -C_y^2    \\
        -C_x^2    & -C_x C_y  & C_x^2     & C_x C_y   \\
        -C_x C_y  & -C_y^2    & C_x C_y   & C_y^2
    \end{bmatrix}
    \begin{matrix}
        \leftarrow u_{iX} \\
        \leftarrow u_{iY} \\
        \leftarrow u_{jX} \\
        \leftarrow u_{jY}
    \end{matrix}
\end{equation}
\end{formulabox}

\textbf{Propiedades Cruciales de $[k_m]$:}
\begin{itemize}
    \item \textbf{Simetría perfecta:} $k_{ij} = k_{ji}$ (Principio de reciprocidad de Maxwell-Betti).
    \item \textbf{Suma nula por filas y columnas:} Cada columna suma cero, garantizando el equilibrio traslacional estático de la barra como cuerpo libre ($\sum F_x = 0, \sum F_y = 0$).
    \item \textbf{Diagonal estrictamente positiva:} $C_x^2 \ge 0$ y $C_y^2 \ge 0$.
\end{itemize}

% -------------------------------------------------------------
\section{Ensamblaje Global: La Matriz de la Estructura (Karmadura)}
% -------------------------------------------------------------

Para una estructura de $N$ nudos, existen $2N$ grados de libertad globales. Para un modelo de hasta 6 nudos, la matriz global $[K_{global}]$ (llamada por el Ing. Ulianov \textit{Karmadura}) tiene una dimensión de:
\begin{equation}
    \text{Dimensión} = 12 \times 12 \quad (U_{1X}, U_{1Y}, U_{2X}, U_{2Y}, \dots, U_{6X}, U_{6Y})
\end{equation}

\subsection{El Principio de Ensamblaje Directo}
Cada barra $m$ aporta su rigidez $4 \times 4$ expandiéndose a una matriz $12 \times 12$ ($[K_m]$) donde los 16 términos se ubican exactamente en las intersecciones de las filas y columnas correspondientes a sus nudos $i$ y $j$:
\begin{equation}
    [K_{global}] = \sum_{m=1}^{B} [K_m]
\end{equation}

\subsection{Significado Físico de la Diagonal Principal}
Los bloques diagonales $2 \times 2$ correspondientes al nudo $n$:
\begin{equation}
    [K_{diag, n}] = \begin{bmatrix} K(U_{nX}, U_{nX}) & K(U_{nX}, U_{nY}) \\ K(U_{nY}, U_{nX}) & K(U_{nY}, U_{nY}) \end{bmatrix}
\end{equation}
representan la \textbf{rigidez propia del nudo $n$}, igual a la suma directa de las rigideces proyectadas de todas las barras conectadas a dicho nudo:
\begin{equation}
    K(U_{nX}, U_{nX}) = \sum_{m \in nudos(n)} \left(\frac{EA}{L}\right)_m C_{x, m}^2
\end{equation}
Los términos fuera de los bloques diagonales representan el \textbf{acoplamiento elástico} entre dos nudos distintos conectados por una barra. Si dos nudos no comparten ninguna barra, su rigidez mutua es idénticamente cero.

% -------------------------------------------------------------
\section{LA GRAN PREGUNTA TEÓRICA: ¿Por qué se eliminan filas y columnas?}
% -------------------------------------------------------------

Esta es la pregunta medular que suele hacer el docente en los exámenes orales y prácticos:
\begin{center}
\textbf{¿Por qué en el método de la rigidez tachamos, eliminamos o particionamos filas y columnas?}
\end{center}

\subsection{La Singularidad Cinemática Inicial}
Antes de aplicar las condiciones de apoyo, la ecuación de equilibrio global es:
\begin{equation}
    \{C\} = [K_{global}] \{D\}
\end{equation}
En este estado inicial, la estructura \textbf{no tiene apoyos vinculados con la tierra}. Por consiguiente, la estructura entera puede trasladarse o rotar libremente en el espacio como un \textbf{cuerpo rígido} sin sufrir deformación interna alguna.

Matemáticamente, esto significa que el determinante de $[K_{global}]$ es \textbf{cero}:
\begin{equation}
    \det([K_{global}]) = 0 \implies [K_{global}] \text{ NO TIENE INVERSA}
\end{equation}
Intentar resolver $\{D\} = [K_{global}]^{-1} \{C\}$ directamente es matemáticamente imposible (división entre cero / matriz singular). Para que la estructura sea estable, debemos imponer las restricciones de apoyo.

\subsection{Partición Cinemática por Condiciones de Apoyo}
Al colocar apoyos en la armadura, el vector total de grados de libertad se divide naturalmente en dos conjuntos disjuntos:
\begin{enumerate}
    \item \textbf{Grados de Libertad Libres (Subíndice 1 o $L$):} Nudos libres o direcciones donde el apoyo permite movimiento (ej. rodillos).
    \begin{itemize}
        \item Desplazamientos: \textbf{INCÓGNITAS DESCONOCIDAS} $\{D_1\}$.
        \item Cargas: \textbf{VALORES CONOCIDOS} $\{C_1\} = \{P_{ext}\}$ (fuerzas actuantes en nudos).
    \end{itemize}
    \item \textbf{Grados de Libertad Restringidos / Apoyos (Subíndice 2 o $R$):} Direcciones bloqueadas por los apoyos con el suelo.
    \begin{itemize}
        \item Desplazamientos: \textbf{VALORES CONOCIDOS Y NULOS} $\{D_2\} = \mathbf{0}$ (condición de apoyo rígido).
        \item Cargas: \textbf{REACCIONES INCÓGNITAS} $\{C_2\} = \{R\}$ (fuerzas que el suelo ejerce para impedir el movimiento).
    \end{itemize}
\end{enumerate}

Reordenando y particionando el sistema matricial global por bloques:
\begin{formulabox}{Ecuación Matricial Particionada Fundamental}
\begin{equation}
    \begin{bmatrix}
        \{C_1\} \\
        \hline
        \{R_2\}
    \end{bmatrix}
    =
    \begin{bmatrix}
        [K_{11}] & [K_{12}] \\
        \hline
        [K_{21}] & [K_{22}]
    \end{bmatrix}
    \begin{bmatrix}
        \{D_1\} \\
        \hline
        \{D_2\}
    \end{bmatrix}
\end{equation}
\end{formulabox}

\subsection{Por qué se eliminan las Filas y Columnas de los Apoyos para hallar Desplazamientos}
Al expandir la primera fila de la ecuación particionada:
\begin{equation}
    \{C_1\} = [K_{11}] \{D_1\} + [K_{12}] \{D_2\}
\end{equation}
Como los desplazamientos en los apoyos son idénticamente nulos ($\{D_2\} = \mathbf{0}$):
\begin{equation}
    [K_{12}] \{D_2\} = [K_{12}] \begin{bmatrix} 0 \\ 0 \\ \vdots \end{bmatrix} = \mathbf{0}
\end{equation}
El término $[K_{12}] \{D_2\}$ se anula por completo. La ecuación queda reducida a:
\begin{equation}
    \{C_1\} = [K_{11}] \{D_1\}
\end{equation}

\begin{conceptbox}{Conclusión Teórica Vital: Justificación de la Eliminación}
\begin{itemize}
    \item \textbf{¿Por qué eliminamos las columnas de los apoyos?} \\
    Porque al multiplicar la matriz de rigidez por el vector de desplazamientos $\{D\}$, las columnas asociadas a los apoyos se multiplican por cero ($D_2 = 0$). No aportan rigidez al cálculo de los desplazamientos libres.
    \item \textbf{¿Por qué eliminamos las filas de los apoyos?} \\
    Porque en las filas de los apoyos las fuerzas $\{C_2\} = \{R_2\}$ son reacciones desconocidas. No podemos resolver una ecuación donde tanto el desplazamiento como la fuerza sean incógnitas simultáneas. Al quedarnos con las filas de los grados libres $\{C_1\}$, tenemos fuerzas externas totalmente conocidas $\{P\}$.
\end{itemize}
\end{conceptbox}

Dado que la estructura posee suficientes apoyos para impedir movimientos de cuerpo rígido, la submatriz $[K_{11}]$ es \textbf{simétrica, definida positiva y estrictamente invertible}:
\begin{formulabox}{Cálculo de Desplazamientos Reales}
\begin{equation}
    \{D_1\} = [K_{11}]^{-1} \{C_1\}
\end{equation}
\end{formulabox}

\subsection{Por qué se utiliza la Submatriz $K_{21}$ para calcular las Reacciones}
Expandiendo la segunda fila del sistema particionado:
\begin{equation}
    \{R_2\} = [K_{21}] \{D_1\} + [K_{22}] \{D_2\}
\end{equation}
Nuevamente, como $\{D_2\} = \mathbf{0}$, el término $[K_{22}] \{D_2\} = \mathbf{0}$. Queda exactamente:
\begin{formulabox}{Cálculo de Reacciones en los Apoyos}
\begin{equation}
    \{R_2\} = [K_{21}] \{D_1\}
\end{equation}
\end{formulabox}
\begin{itemize}
    \item $[K_{21}]$ es la \textbf{submatriz de acoplamiento elástico}: contiene las filas de los apoyos y las columnas de los grados libres.
    \item Físicamente, representa cómo las deformaciones de los nudos libres empujan y transmiten carga a los apoyos fijos o rodillos.
\end{itemize}

% -------------------------------------------------------------
\section{Cálculo de Fuerzas Axiales en las Barras y Diagnóstico de Estado}
% -------------------------------------------------------------

Una vez calculados todos los desplazamientos nodales globales $\{D\}$, para cada barra $m$ que conecta el nudo $i$ con el nudo $j$:
\begin{equation}
    \{D_m\} = \begin{bmatrix} D_{iX} \\ D_{iY} \\ D_{jX} \\ D_{jY} \end{bmatrix}
\end{equation}

La fuerza axial interna $F_m$ viene dada por:
\begin{formulabox}{Fórmula de Fuerza Axial (Ing. Ulianov)}
\begin{equation}
    F_m = \left(\frac{EA}{L}\right)_m \begin{bmatrix} -C_x & -C_y & C_x & C_y \end{bmatrix} \begin{bmatrix} D_{iX} \\ D_{iY} \\ D_{jX} \\ D_{jY} \end{bmatrix}
\end{equation}
\end{formulabox}

\subsection{Criterio de Signos y Diagnóstico de Estado}
\begin{itemize}
    \item Si $F_m > 0 \implies \mathbf{BARRA\ A\ TRACCI\acute{O}N\ (+)}$: La barra experimenta tensión pura; tiende a alargarse.
    \item Si $F_m < 0 \implies \mathbf{BARRA\ A\ COMPRESI\acute{O}N\ (-)}$: La barra experimenta empuje axial; tiende a acortarse (debe verificarse contra pandeo por Euler).
    \item Si $F_m = 0 \implies \mathbf{BARRA\ DE\ FUERZA\ NULA}$: Barra de estabilidad o arriostramiento sin carga axial bajo el estado actual.
\end{itemize}

\subsection{Verificación de Equilibrio Estático Global}
Para cualquier armadura en equilibrio, las reacciones obtenidas deben satisfacer idénticamente:
\begin{equation}
    \sum F_X = \sum P_X + \sum R_X = 0, \quad \sum F_Y = \sum P_Y + \sum R_Y = 0, \quad \sum M_O = 0
\end{equation}

% -------------------------------------------------------------
\section{Guía Tutorial: Cómo Construir la Plantilla en Excel Paso a Paso}
% -------------------------------------------------------------

Si el docente te pide en clase o en un examen: \textit{"Construye la hoja desde cero en un libro en blanco"}, este es el procedimiento exacto paso a paso:

\subsection{Paso 1: Tabla de Coordenadas, Apoyos y Cargas (Columnas P a W)}
\begin{enumerate}
    \item En las columnas P a W (filas 5 a 11), crea los encabezados:
    \texttt{NUDO | X (cm) | Y (cm) | TIPO APOYO | Rx | Ry | Px (Kgf) | Py (Kgf)}.
    \item En la columna \texttt{TIPO APOYO}, ingresa: \texttt{Fijo}, \texttt{Móvil Y}, \texttt{Móvil X} o \texttt{Libre}.
    \item En la columna $Rx$ (Col T, fila 6), coloca la fórmula que detecta apoyo en X:
    \begin{verbatim}
    =IF(Q6="","", IF(LEFT(UPPER(TRIM(S6)),4)="FIJO", 1,
     IF(OR(UPPER(TRIM(S6))="MOVIL X", UPPER(TRIM(S6))="MÓVIL X",
           UPPER(TRIM(S6))="RODILLO X"), 1, 0)))
    \end{verbatim}
    \item En la columna $Ry$ (Col U, fila 6), coloca la fórmula que detecta apoyo en Y:
    \begin{verbatim}
    =IF(Q6="","", IF(LEFT(UPPER(TRIM(S6)),4)="FIJO", 1,
     IF(OR(UPPER(TRIM(S6))="MOVIL Y", UPPER(TRIM(S6))="MÓVIL Y",
           UPPER(TRIM(S6))="RODILLO Y", UPPER(TRIM(S6))="MOVIL"), 1, 0)))
    \end{verbatim}
\end{enumerate}

\subsection{Paso 2: Conectividad y Grados de Libertad (Tabla 2, Filas 17 a 24)}
\begin{enumerate}
    \item Para cada barra (filas 17 a 24):
    \begin{itemize}
        \item Columna C: Nudo Inicial (ej. 1).
        \item Columna D: \texttt{=IF(C17="","", "U" \& C17 \& "X")}.
        \item Columna E: \texttt{=IF(C17="","", "U" \& C17 \& "Y")}.
        \item Columna F: Nudo Final (ej. 3).
        \item Columna G: \texttt{=IF(F17="","", "U" \& F17 \& "X")}.
        \item Columna H: \texttt{=IF(F17="","", "U" \& F17 \& "Y")}.
    \end{itemize}
\end{enumerate}

\subsection{Paso 3: Propiedades Geométricas de las Barras (Tabla 1, Filas 6 a 13)}
\begin{enumerate}
    \item Coordenadas dinámicas mediante \texttt{VLOOKUP}:
    \begin{itemize}
        \item $X_{ini}$: \texttt{=IF(C17="","", VLOOKUP(C17, \$P\$6:\$W\$11, 2, FALSE))}.
        \item $Y_{ini}$: \texttt{=IF(C17="","", VLOOKUP(C17, \$P\$6:\$W\$11, 3, FALSE))}.
        \item $X_{fin}$: \texttt{=IF(F17="","", VLOOKUP(F17, \$P\$6:\$W\$11, 2, FALSE))}.
        \item $Y_{fin}$: \texttt{=IF(F17="","", VLOOKUP(F17, \$P\$6:\$W\$11, 3, FALSE))}.
    \end{itemize}
    \item $DX = X_{fin} - X_{ini}$, $DY = Y_{fin} - Y_{ini}$.
    \item $L = \text{SQRT}(DX^2 + DY^2)$.
    \item $Cx = DX / L$, $Cy = DY / L$.
    \item Rigidez axial: $AE/L = (A \cdot E) / L$.
\end{enumerate}

\subsection{Paso 4: Vectores Dinámicos \{C\} y \{D\} y Columnas Auxiliares F y G}
Para cada grado de libertad $r$:
\begin{itemize}
    \item \textbf{Vector C (Cargas/Reacciones):}
    \begin{verbatim}
    =IF(Q_coord="", 0, IF(Rx=1, B_var, Px))
    \end{verbatim}
    Si hay apoyo ($Rx=1$), coloca el nombre de la reacción incógnita (\texttt{"C1X"}). Si está libre ($Rx=0$), coloca la carga nodal conocida (\texttt{Px}).
    \item \textbf{Vector D (Desplazamientos):}
    \begin{verbatim}
    =IF(Q_coord="", 0, IF(Rx=1, 0, B_var))
    \end{verbatim}
    Si hay apoyo ($Rx=1$), el desplazamiento es \texttt{0}. Si está libre ($Rx=0$), es la incógnita (\texttt{"D1X"}).
    \item \textbf{Columna Auxiliar F (Conteo de Grados Libres):}
    \begin{verbatim}
    F_primera_fila: =IF(ISTEXT(E_disp), 1, "")
    F_siguiente:    =IF(ISTEXT(E_disp), MAX(F$start:F_prev) + 1, "")
    \end{verbatim}
    Enumera correlativamente $1, 2, \dots, n_{free}$ solo los grados libres reales.
    \item \textbf{Columna Auxiliar G (Conteo de Reacciones):}
    \begin{verbatim}
    G_primera_fila: =IF(ISTEXT(E_carga), 1, "")
    G_siguiente:    =IF(ISTEXT(E_carga), MAX(G$start:G_prev) + 1, "")
    \end{verbatim}
    Enumera correlativamente $1, 2, \dots, n_{rest}$ solo los apoyos reales.
\end{itemize}

\subsection{Paso 5: Extracción de $K_{11}$ e Inversión Regularizada (Cero \#¡NUM!)}
\begin{enumerate}
    \item Los encabezados de fila y columna se extraen con \texttt{INDEX/MATCH} buscando el índice $1, 2, \dots$ en la columna F.
    \item \textbf{Fórmula Regularizada en cada celda $K_{11}(i, j)$:}
    \begin{verbatim}
    =IF(OR(row_dof="-", col_dof="-"), IF(i=j, 1, 0),
     IF(OR(diag_row=0, diag_col=0), IF(i=j, 1, 0),
     INDEX(Karmadura, MATCH(row_dof), MATCH(col_dof))))
    \end{verbatim}
    \item En la matriz inversa:
    \begin{verbatim}
    {=MINVERSA(K11_range)}
    \end{verbatim}
    \item En los desplazamientos libres:
    \begin{verbatim}
    {=MMULT(MINVERSA_range, C_free_range)}
    \end{verbatim}
\end{enumerate}

\subsection{Paso 6: Cálculo de Reacciones con $K_{21}$}
\begin{enumerate}
    \item Las filas de $K_{21}$ se extraen buscando el índice $1, 2, \dots$ en la columna G (reacciones).
    \item Las columnas son los grados libres (mismas columnas que $K_{11}$).
    \item Reacciones finales:
    \begin{verbatim}
    {=MMULT(K21_range, D_free_range)}
    \end{verbatim}
\end{enumerate}

% -------------------------------------------------------------
\section{Ejemplo Numérico Resuelto: Ejercicio del Alumno (3 Nudos, 3 Barras)}
% -------------------------------------------------------------

\subsection{Datos de Entrada}
\begin{itemize}
    \item \textbf{Nudo 1:} $X_1 = 0$, $Y_1 = 0$. Apoyo \textbf{Móvil Y} ($R_x=0, R_y=1$). Carga $P = (0, 0)$.
    \item \textbf{Nudo 2:} $X_2 = -700\text{ cm}$, $Y_2 = 0$. Apoyo \textbf{Fijo} ($R_x=1, R_y=1$). Carga $P = (0, 0)$.
    \item \textbf{Nudo 3:} $X_3 = -400\text{ cm}$, $Y_3 = 500\text{ cm}$. Nudo \textbf{Libre} ($R_x=0, R_y=0$). Carga $P = (4000\text{ Kgf}, -5000\text{ Kgf})$.
    \item Material y sección: $A = 10\text{ cm}^2$, $E = 2,100,000\text{ Kgf/cm}^2$.
    \item Conectividad: Barra 1 ($1 \to 3$), Barra 2 ($2 \to 3$), Barra 3 ($2 \to 1$).
\end{itemize}

\subsection{Grados de Libertad Resultantes}
\begin{itemize}
    \item \textbf{Grados Libres ($n_{free} = 3$):} $U_{1X}$, $U_{3X}$, $U_{3Y}$.
    \item \textbf{Grados Restringidos ($n_{rest} = 3$):} $U_{1Y}$ ($R_{1Y}$), $U_{2X}$ ($R_{2X}$), $U_{2Y}$ ($R_{2Y}$).
\end{itemize}

\subsection{Resultados Numéricos Finales}
\begin{table}[h]
\centering
\begin{tabular}{lll}
\toprule
\textbf{Variable} & \textbf{Valor Numérico} & \textbf{Interpretación Física} \\
\midrule
$D_{1X}$ & $+0.13333333\text{ cm}$ & El rodillo se desplaza hacia la derecha \\
$D_{1Y}$ & $0.00000000\text{ cm}$ & Apoyo móvil restringe desplazamiento vertical \\
$D_{2X}$ & $0.00000000\text{ cm}$ & Apoyo fijo articulado: ¡Cero desplazamiento! \\
$D_{2Y}$ & $0.00000000\text{ cm}$ & Apoyo fijo articulado: ¡Cero desplazamiento! \\
$D_{3X}$ & $+0.25478102\text{ cm}$ & Desplazamiento horizontal del nudo libre \\
$D_{3Y}$ & $-0.15286861\text{ cm}$ & Desplazamiento vertical hacia abajo por carga \\
\midrule
$R_{1Y}$ & $+5000.00\text{ Kgf}$ & Única reacción en Nudo 1 (hacia arriba) \\
$R_{2X}$ & $-4000.00\text{ Kgf}$ & Reacción horizontal en Nudo 2 (hacia la izquierda) \\
$R_{2Y}$ & $0.00\text{ Kgf}$ & Reacción vertical nula por equilibrio \\
\midrule
Barra 1 ($1 \to 3$) & $-6403.12\text{ Kgf}$ & \textbf{BARRA A COMPRESIÓN (-)} \\
Barra 2 ($2 \to 3$) & $0.00\text{ Kgf}$ & \textbf{BARRA DE FUERZA NULA} \\
Barra 3 ($2 \to 1$) & $+4000.00\text{ Kgf}$ & \textbf{BARRA A TRACCIÓN (+)} \\
\bottomrule
\end{tabular}
\caption{Resultados exactos validados del ejercicio resuelto.}
\end{table}

\begin{examtip}{Pregunta de Examen Típica del Ing. Ulianov}
\textit{"¿Por qué la Barra 2 tiene fuerza axial nula ($N_2 = 0.00$) si está conectada a la carga de 4000 y -5000 en el Nudo 3?"} \\
\textbf{Respuesta Maestra:} \\
Porque la barra 2 está orientada exactamente a lo largo de la resultante que no requiere resistir corte en el nudo, y al resolver el equilibrio nodal en el nudo 3 o por compatibilidad cinemática, la deformación axial $\Delta L_2 = [-C_{x2}, -C_{y2}, C_{x2}, C_{y2}] \cdot \{D\}$ resulta exactamente cero ($0.00000000$).
\end{examtip}

\end{document}
